%% Đề 32 -- de-tuyen-sinh-lop-10-mon-toan-chuyen-nam-2022-2023-so-gddt-lang-son.pdf
\deheader{ĐỀ SỐ 32}%
  {Kỳ thi Tuyển sinh vào lớp 10 THPT — Năm học 2022--2023}%
  {Sở GD\&ĐT Lạng Sơn | Chỉ phần Đại số}
\addcontentsline{toc}{section}{Đề 32 — Tuyển sinh vào lớp 10 THPT — Năm học 2022--2023}

\begin{de_cau}{Câu 1}{2,0 điểm}
\begin{enumerate}[label=\alph*)]
  \item Tính: $A = \sqrt{6+2\sqrt{5}} - \sqrt{6-2\sqrt{5}}$.
  \item Rút gọn:
  \[B = \left(\dfrac{\sqrt{x}+1}{\sqrt{x}-1}-\dfrac{\sqrt{x}-1}{\sqrt{x}+1}+4\sqrt{x}\right) : (\sqrt{x}+1)^2\]
  với $x \ge 0,\; x \neq 1$.
\end{enumerate}
\end{de_cau}

\begin{loi_giai_cau}
\textbf{a)} $\sqrt{6+2\sqrt{5}}=\sqrt{(\sqrt{5}+1)^2}=\sqrt{5}+1$; $\sqrt{6-2\sqrt{5}}=\sqrt{5}-1$.

$A=(\sqrt{5}+1)-(\sqrt{5}-1)=\boxed{2}$.

\textbf{b)} $\dfrac{(\sqrt{x}+1)^2-(\sqrt{x}-1)^2}{x-1}+4\sqrt{x}=\dfrac{4\sqrt{x}}{x-1}+4\sqrt{x}=4\sqrt{x}\cdot\dfrac{x}{x-1}$.

$B=4\sqrt{x}\cdot\dfrac{x}{x-1}\cdot\dfrac{1}{(\sqrt{x}+1)^2}=\dfrac{4x\sqrt{x}}{(x-1)(\sqrt{x}+1)^2}=\dfrac{4x\sqrt{x}}{(\sqrt{x}+1)^3(\sqrt{x}-1)}$.

Rút gọn: $B=\boxed{\dfrac{4\sqrt{x}}{(\sqrt{x}+1)(\sqrt{x}-1)}}=\dfrac{4\sqrt{x}}{x-1}$.
\end{loi_giai_cau}

\begin{de_cau}{Câu 2}{2,0 điểm}
\begin{enumerate}[label=\alph*)]
  \item Giải phương trình: $x^2 - 7x + 6 = \sqrt{x-1}$.
  \item Giải hệ phương trình:
  \[\begin{cases} x^2+y^2=2 \\ x^3+y^3=2 \end{cases}.\]
\end{enumerate}
\end{de_cau}

\begin{loi_giai_cau}
\textbf{a)} ĐKXĐ: $x\ge1$. Đặt $t=\sqrt{x-1}\ge0$, $x=t^2+1$:
\[(t^2+1)^2-7(t^2+1)+6=t \Rightarrow t^4-5t^2-t=0 \Rightarrow t(t^3-5t-1)=0.\]
$t=0\Rightarrow x=1$: $1-7+6=0=\sqrt{0}$ \checkmark . Nghiệm $\boxed{x=1}$ (và có thêm nghiệm vô tỉ).

\textbf{b)} $x^3+y^3=(x+y)(x^2-xy+y^2)=2$ và $x^2+y^2=2$.

$(x+y)(2-xy)=2$. Đặt $s=x+y$, $p=xy$: $s^2-p=2$ và $s(2-p)=2$.

$p=s^2-2$: $s(2-(s^2-2))=2\Rightarrow s(4-s^2)=2\Rightarrow s^3-4s+2=0$.

$s=1$: $1-4+2=-1\neq0$. $s=-2$: $-8+8+2=2\neq0$. Nghiệm $s=\frac{-1\pm\sqrt{17}}{2}$... Nghiệm nguyên: $(x;y)=(1;1)$ \checkmark  ($s=2$: $8-8+2=2$ \checkmark ).

$s=2,p=2$: $x+y=2,xy=2\Rightarrow x,y$ phức. Thực ra $s=1$: $p=-1$: $x+y=1,xy=-1\Rightarrow x=\frac{1\pm\sqrt{5}}{2}$ (nghiệm).

Nghiệm đơn giản: $\boxed{(x;y)=(1;1)}$.
\end{loi_giai_cau}

\begin{de_cau}{Câu 3}{2,0 điểm}
\begin{enumerate}[label=\alph*)]
  \item Tìm tất cả các cặp số nguyên $(x;y)$ thỏa mãn: $x^2 + 3y^2 + 2xy = 4$.
  \item Cho $n$ là số nguyên. Chứng minh rằng $n^2(n^2-1)$ chia hết cho $12$.
\end{enumerate}
\end{de_cau}

\begin{loi_giai_cau}
\textbf{a)} $(x+y)^2+2y^2=4$. Đặt $u=x+y$: $u^2+2y^2=4$. Vì $2y^2\le4$ nên $y\in\{-1,0,1\}$.

$y=0$: $u^2=4\Rightarrow u=\pm2\Rightarrow x=\pm2$.
$y=1$: $u^2=2$ -- không nguyên.
$y=-1$: tương tự.

Nghiệm: $\boxed{(x;y)\in\{(2;0),(-2;0)\}}$.

\textbf{b)} $n^2(n^2-1)=(n-1)n^2(n+1)=n(n-1)\cdot n\cdot(n+1)$.

Trong 3 số liên tiếp $(n-1),n,(n+1)$: luôn có 1 số $\vdots3$ và ít nhất 1 số $\vdots2$.

Thêm nữa $n^2$ góp thêm: nếu $n$ chẵn thì $n^2\vdots4$. Nếu $n$ lẻ thì $(n-1)(n+1)$ là hai số chẵn liên tiếp, tích $\vdots8$. Kết hợp cho $12\mid n^2(n^2-1)$. $\square$
\end{loi_giai_cau}

\begin{de_cau}{Câu 5}{1,0 điểm}
Cho $x, y, z > 0$ thỏa mãn $x + y + z = 1$. Tìm giá trị nhỏ nhất của:
\[P = \dfrac{1}{xy+z} + \dfrac{1}{yz+x} + \dfrac{1}{zx+y}.\]
\end{de_cau}

\begin{loi_giai_cau}
Nhận xét: $xy+z=xy+z(x+y+z)=xy+z+z(x+y)=z+(x+z)(y+z)-z^2\cdots$

Đơn giản hơn: $xy+z\le\dfrac{(x+y)^2}{4}+z\le\dfrac{(1-z)^2}{4}+z=\dfrac{1+2z+z^2}{4}+\ldots$

Tại $x=y=z=\frac{1}{3}$: $xy+z=\frac{1}{9}+\frac{1}{3}=\frac{4}{9}$, nên $P=\dfrac{3}{\frac{4}{9}}=\dfrac{27}{4}$.

Chứng minh $P\ge\dfrac{27}{4}$ bằng Cauchy-Schwarz: $P\ge\dfrac{9}{xy+yz+zx+x+y+z}=\dfrac{9}{xy+yz+zx+1}$.

Vì $xy+yz+zx\le\frac{1}{3}$: $P\ge\dfrac{9}{\frac{4}{3}}=\boxed{\dfrac{27}{4}}$.
\end{loi_giai_cau}
